% LỜI TỰA --- Quyển XXI
\chapter*{Lời Tựa}
\addcontentsline{toc}{chapter}{Lời Tựa}
\markboth{Lời Tựa}{Lời Tựa}

\begin{truyen}{Lời Tựa}{30 Wake-Up Calls for the Reluctant Student}
\chuhoa{G}{iáo viên bắt đầu bộ sách này với một câu hỏi đơn giản: làm thế nào để mỗi buổi học có thể để lại thứ gì đó đáng nhớ hơn là bài thi?}
\textit{(A teacher began this series with a simple question: how can each class leave behind something more memorable than a test?)}

Có những câu thầy nói không la mắng, không dọa~--- nhưng học sinh lười nghe xong giật mình.
\textit{(Some things the teacher says are not scolding or threats~--- but lazy students startle when they hear them.)}

Quyển hai mươi mốt gom 30 câu như thế: thẳng, rõ, để học sinh tỉnh ra trước khi trả giá bằng điểm kém, bằng cơ hội mất, bằng hối tiếc sau này. Mong thầy cô dùng đúng lúc; mong học trò đọc khi còn kịp.
\textit{(Volume twenty-one gathers 30 such sentences: direct and clear, so students wake up before paying with bad grades, lost chances, and future regret. May teachers use them at the right moment; may students read them while there is still time.)}

Bộ sách <<Truyện Tích Cảm Hứng Song Ngữ>> gồm 24 quyển, mỗi quyển một chủ đề riêng. Quyển XXI này mang chủ đề: \textbf{thẳng, rõ, để tỉnh ra trước khi trả giá}.
\textit{(The series ``Inspiring Bilingual Tales'' contains 24 volumes, each with its own theme. Volume XXI focuses on: \textbf{thẳng, rõ, để tỉnh ra trước khi trả giá}.})
\end{truyen}

\begin{baihoc}
Giật mình một lần có thể thay đổi hướng đi~--- nếu em chịu lắng nghe và hành động.
\textit{(One wake-up moment can change your direction~--- if you listen and act.)}
\end{baihoc}

\begin{ghinhoanh}
\textbf{Short English to remember:}

A wake-up call now saves a lifetime of regret.

The cost of laziness is always paid with interest.
\end{ghinhoanh}

\begin{tuhocnhanh}
1. Câu nào trong 30 câu này khiến em giật mình nhất? Vì sao? \textit{(Which of the 30 sentences gave you the biggest wake-up? Why?)}

2. Em đang trì hoãn điều gì mà em biết mình sẽ hối tiếc nếu không làm sớm? \textit{(What are you delaying that you know you'll regret if you don't act soon?)}

3. Điều gì giúp em từ lười trở nên chịu cố gắng? \textit{(What helps you go from lazy to actually trying?)}
\end{tuhocnhanh}

\begin{center}
\textit{\color{accentsgold}``Lười hôm nay, trả giá ngày mai.''}
\end{center}

\begin{flushright}
\textbf{Tôi Nguyễn Văn Sang}\\
\textit{GV THPT Nguyễn Hữu Cảnh - TP HCM}
\end{flushright}

\ngancach
